\documentclass[11pt]{article}
\usepackage[letterpaper,margin=1in]{geometry}
\usepackage{amsmath,amssymb,amsthm}
\usepackage{fontspec}
\setmainfont{lmroman10-regular.otf}[BoldFont=lmroman10-bold.otf,ItalicFont=lmroman10-italic.otf,BoldItalicFont=lmroman10-bolditalic.otf]
\setsansfont{lmsans10-regular.otf}
\setmonofont{lmmono10-regular.otf}
\usepackage{microtype}
\usepackage{booktabs}
\makeatletter
\renewcommand\section{\@startsection{section}{1}{\z@}
  {-3.5ex\@plus -1ex\@minus -.2ex}{2.3ex\@plus .2ex}
  {\normalfont\Large\bfseries\raggedright}}
\makeatother
\usepackage{xurl}
\usepackage[unicode,hidelinks,hypertexnames=false]{hyperref}
\hypersetup{pdftitle={Ternary derivative controls and Reed-Muller weight caps: C44 < 2.1974348818},pdfauthor={Ricardo Pinho},pdfsubject={Candidate, unrefereed; frozen proof with explicit corrections}}
\newtheorem*{claim}{Statement}
\theoremstyle{remark}
\newtheorem*{editorial}{Editorial remark}
\newtheoremstyle{correctionstyle}{0.6em}{0.6em}{\normalfont}{}{\bfseries}{:}{0.5em}{}
\theoremstyle{correctionstyle}
\newtheorem*{correction}{Correction}
\setlength{\parskip}{0.25em}
\setlength{\parindent}{0pt}
\setlength{\emergencystretch}{2em}
\allowdisplaybreaks[1]
\title{Ternary derivative controls and Reed-Muller weight caps: C44 < 2.1974348818}
\author{Ricardo Pinho\\\small ORCID \href{https://orcid.org/0009-0006-2744-2599}{0009-0006-2744-2599}}
\date{4 November 2026; Section 4 opening expanded 6 November 2026}
\begin{document}
\maketitle
\begin{center}
\fbox{\begin{minipage}{0.94\linewidth}
Candidate, unrefereed. This result was found by AI agents working under my direction. I have only reviewed it at summary level, and other AI models reviewed it, so no human has checked it line by line.

The verifiers check the arithmetic, not the written lemmas.
\end{minipage}}
\end{center}
\begin{claim}
For the constant in $r(f)\le C\,2^{\deg(f)}$,
\[
C_{44}\le\frac{174098727923678791596986402933}{79228162514264337593543950336}<2.1974348818.
\]
The lower bound is $1.5$.
\end{claim}
\begin{editorial}
The comparison with the exact rational is non-strict, as in the frozen
proof's equation (12). The source README states it as strict; no strictness
argument at the rational endpoint is supplied here. Since the rational is
strictly below $2.1974348818$, the decimal bound $C_{44}<2.1974348818$ holds either way.
\end{editorial}
The seven numbered sections below transcribe the frozen \path{PROOF.md}.
The exception is Section 4, which was expanded step by step and reworded for
clarity on 6 November 2026, as marked there.
Marked corrections reproduce the supplied errata without silently changing
the original argument. Supplemental material follows Section 7.

% SOURCE P001, PROOF.md lines 1--1
\textit{Frozen source title: Ternary derivative controls and higher weight caps}\phantomsection\label{src-p001}\label{src-end-p001}

% SOURCE P002, PROOF.md lines 3--8
\leavevmode\label{src-p002}Candidate, unrefereed. This is a written mathematical proof with exact
arithmetic certificates, not proof-assistant formalization or independent
human refereeing. Two separately written implementations check the finite
artifact. Their finite controls do not prove the named external theorems.
No novelty or priority is claimed. The framework is due to CHS/Wellens; the
low binary weights are due to Berlekamp–Sloane and Kasami–Tokura.\label{src-end-p002}\par

% SOURCE P003, PROOF.md lines 10--10
\clearpage
\section{Definitions, support, and extremizers}\label{src-p003}\label{src-end-p003}

% SOURCE P004, PROOF.md lines 12--16
\leavevmode\label{src-p004}The cube has independent uniform sign coordinates. Every real function has
a unique multilinear Fourier expansion $f=\sum_S\widehat f(S)\prod_{i\in S}x_i$.
Its real degree is the largest $|S|$ with nonzero coefficient. A coordinate
is relevant precisely when some such $S$ contains it. Its local degree $d_i$
is the largest $|S|$ containing it. Set\label{src-end-p004}\par

% SOURCE P005, PROOF.md lines 18--20
\begin{equation}
\begin{aligned}
D(f)&=\sum_{i\ \mathrm{relevant}}2^{-d_i},\\
R_d&=\max\text{ number of relevant coordinates at degree at most }d,\\
W_d&=\max D(f)\text{ at degree at most }d.
\end{aligned}\tag*{}
\label{src-p005}
\end{equation}

% SOURCE P006, PROOF.md lines 22--29
\leavevmode\label{src-p006}A nonzero multilinear real polynomial of degree at most $k$ is nonzero on
at least a $2^{-k}$ fraction of the cube. Here is an induction proof. The case
$k=0$ is constant. For a nonconstant polynomial choose a coordinate $j$ and
write $P=x_jA+B$ with $A$ nonzero and degree at most $k-1$. The induction
hypothesis gives support density of $A$ at least $2^{-(k-1)}$. For each input
in that support, the two values $B+A$ and $B-A$ cannot both vanish. Averaging
the two $j$ branches gives support density at least $2^{-k}$. This also handles
polynomials whose degree is strictly below the stated upper bound.\label{src-end-p006}\par

% SOURCE P007, PROOF.md lines 31--31
\leavevmode\label{src-p007}For a sign Boolean $f$, write\label{src-end-p007}\par

% SOURCE P008, PROOF.md lines 33--33
\begin{equation}
A_i=\frac{f(i=+1)-f(i=-1)}{2}.\tag*{}
\label{src-p008}
\end{equation}

% SOURCE P009, PROOF.md lines 35--36
\leavevmode\label{src-p009}Its values are $0,+1,-1$, its real degree is $d_i-1$, and $p_i=\mathbb E A_i^2=\operatorname{Inf}_i(f)$.
The support lemma gives $p_i\ge 2^{1-d_i}$. Parseval gives\label{src-end-p009}\par

% SOURCE P010, PROOF.md lines 38--38
\begin{equation}
I(f)=\sum_i p_i=\sum_S |S|\widehat f(S)^2\le\deg(f).\tag*{}
\label{src-p010}
\end{equation}

% SOURCE P011, PROOF.md lines 40--43
\leavevmode\label{src-p011}Consequently $D(f)\le I(f)/2\le d/2$ and every relevant variable has influence
at least $2^{1-d}$. Thus $R_d\le d\,2^{d-1}$. After removing dummy variables and
relabeling, only finitely many Boolean functions need be considered. The
two displayed maxima therefore exist.\label{src-end-p011}\par

% SOURCE P012, PROOF.md lines 45--49
\leavevmode\label{src-p012}We use the CHS identity $W_d=R_d/2^d$, and include its proof because its
extremal reduction matters here. Choose $f$ maximizing $D$ at degree at most $d$.
Let $L$ be its relevant coordinates of local degree below $d$, and $Y$ the
others. If $L$ is nonempty, take a disjoint copy $L'$ and a fresh sign selector $z$,
and form\label{src-end-p012}\par

% SOURCE P013, PROOF.md lines 51--51
\begin{equation}
h=\frac{1+z}{2}f(x_L,Y)+\frac{1-z}{2}f(x_{L'},Y).\tag*{}
\label{src-p013}
\end{equation}

% SOURCE P014, PROOF.md lines 53--62
\leavevmode\label{src-p014}This is Boolean. All degree-$d$ terms of the two $f$ copies involve only $Y$,
so those terms cancel in the coefficient of $z$. Hence $\deg(h)\le d$. Every $Y$
coordinate retains local degree $d$. Each original $L$ coordinate and its
copy have local degree exactly one more than before: a largest Fourier
term containing that coordinate gives a noncanceling term also containing $z$.
Their two weights sum to the original weight. The selector is relevant
because the two copies differ for some inputs whenever $L$ is nonempty.
It adds positive weight. Thus $D(h)>D(f)$, a contradiction. An extremizer
for $D$ has every relevant local degree $d$, proving $W_d\le R_d/2^d$. Conversely,
an $R_d$ extremizer has $D\ge R_d/2^d$. This proves the identity.\label{src-end-p014}\par

% SOURCE P015, PROOF.md lines 64--73
\leavevmode\label{src-p015}We also need every coordinate of an $R_d$ extremizer to have local degree $d$.
In bit coordinates, replace a coordinate of smaller local degree by the
AND of two fresh bits. The new polynomial has degree at most $d$, since
terms using that coordinate increase degree by at most one, and other
terms keep their degrees. Both fresh coordinates are relevant: fix the
other fresh bit to 1 and use a witnessing pair of inputs for the original
coordinate. All other relevant coordinates remain relevant, by realizing
either value of the substituted AND. The relevant count increases by one,
contradicting maximality. This proves the required all-local-degrees-$d$
reduction for count extremizers separately.\label{src-end-p015}\par

% SOURCE P016, PROOF.md lines 75--83
\leavevmode\label{src-p016}The published exact result $R_3=10$ supplies $W_3=5/4$. The elementary influence
bound supplies $W_1\le 1/2$ and $W_2\le 1$. The $R_3$ premise is the result for
arbitrary degree-at-most-three Boolean functions, not a restricted coloring
class; see Tarannikov–Kirienko, ePrint 2000/050, and Krotov–Valyuzhenich,
\url{https://arxiv.org/html/2311.05566v2} , section 3. CHS's identity is Corollary
2.2 in the final Combinatorica paper, \url{https://par.nsf.gov/servlets/purl/10201042} ;
the corrected duplication proof also appears in Chiarelli's thesis,
Proposition 8.5 and Corollary 8.6,
\url{https://sites.math.rutgers.edu/~zeilberg/Theses/JohnChiarelli.pdf} .\label{src-end-p016}\par

% SOURCE P017, PROOF.md lines 85--85
\section{A general control bound for ternary polynomials}\label{src-p017}\label{src-end-p017}

% SOURCE P018, PROOF.md lines 87--88
\leavevmode\label{src-p018}Let $A$ take values $0,+1,-1$, have real degree at most $k\ge 1$, and have $h$
relevant coordinates. Put $p=\mathbb E A^2$. For a relevant coordinate $j$,\label{src-end-p018}\par

% SOURCE P019, PROOF.md lines 90--90
\begin{equation}
\Delta_jA=\frac{A(j=+1)-A(j=-1)}{2}\tag*{}
\label{src-p019}
\end{equation}

% SOURCE P020, PROOF.md lines 92--94
\leavevmode\label{src-p020}is a nonzero real polynomial of degree at most $k-1$. Its nonzero values
belong to $\{+1/2,\allowbreak-1/2,\allowbreak+1,\allowbreak-1\}$, so their squared magnitude is at least $1/4$.
By the support lemma,\label{src-end-p020}\par

% SOURCE P021, PROOF.md lines 96--96
\begin{equation}
\mathbb E(\Delta_jA)^2\ge\frac{1}{4}\,2^{-(k-1)}=2^{-(k+1)}.\tag*{}
\label{src-p021}
\end{equation}

% SOURCE P022, PROOF.md lines 98--99
\leavevmode\label{src-p022}On the other hand Fourier orthogonality, first for each derivative and
then summed over coordinates, gives\label{src-end-p022}\par

% SOURCE P023, PROOF.md lines 101--102
\begin{equation}
\sum_j\mathbb E(\Delta_jA)^2=\sum_S|S|\widehat A(S)^2\le k\sum_S\widehat A(S)^2=kp.\tag*{}
\label{src-p023}
\end{equation}

% SOURCE P024, PROOF.md lines 104--104
\leavevmode\label{src-p024}There are $h$ nonzero summands, each at least $2^{-(k+1)}$. Therefore\label{src-end-p024}\par

% SOURCE P025, PROOF.md lines 106--106
\begin{equation}
h\le\left\lfloor kp\,2^{k+1}\right\rfloor.\tag{1}
\label{src-p025}
\end{equation}

% SOURCE P026, PROOF.md lines 108--109
\leavevmode\label{src-p026}For the derivative $A_i$ of a $\text{degree-}d$ Boolean function set $k=d-1$,
$M=2^k$ and $\lambda=Mp_i$. Formula (1) becomes\label{src-end-p026}\par

% SOURCE P027, PROOF.md lines 111--111
\begin{equation}
h_i\le\lfloor2k\lambda\rfloor.\tag{2}
\label{src-p027}
\end{equation}

% SOURCE P028, PROOF.md lines 113--118
\leavevmode\label{src-p028}No affine change of coordinates over $\mathbb F_2$ is being used to infer real
polynomial degree. Formula (1) is a direct real Fourier calculation. At
$\lambda=1$ it recovers the worst $2k$ budget in the minimum-support
classification of Valyuzhenich, Theorem 2(2),
\url{https://arxiv.org/html/2303.10995} . That external classification is useful
context but is not needed as a premise of (1).\label{src-end-p028}\par

% SOURCE P029, PROOF.md lines 120--120
\section{Restrictions controlled by any nonzero derivative}\label{src-p029}\label{src-end-p029}

% SOURCE P030, PROOF.md lines 122--123
\leavevmode\label{src-p030}For completeness, the weighted restriction inequality for a coordinate
set $H$ is\label{src-end-p030}\par

% SOURCE P031, PROOF.md lines 125--125
\begin{equation}
D(f)\le\sum_{j\in H}2^{-d_j(f)}+\mathbb E_\alpha D(f_\alpha).\tag{3}
\label{src-p031}
\end{equation}

% SOURCE P032, PROOF.md lines 127--137
\leavevmode\label{src-p032}To prove it for one fixed coordinate $j$, consider another relevant
coordinate $i$. If $i$ survives in both branches, its local degrees cannot
increase, so both branch weights are at least its original weight. If it
survives in exactly one branch, its derivative factors as $(1\pm x_j)/2$ times
the nonzero branch derivative. Its original local degree is exactly one
larger than its surviving branch local degree. Half that branch weight
therefore equals the original weight. A relevant $i$ cannot disappear in
both branches. Sum these inequalities over $i$ other than $j$. Applying them
successively to coordinates in $H$, or equivalently inducting on $|H|$, proves
(3) for the coordinates outside $H$; the original weights of $H$ supply the
first term. Disappearing coordinates contribute zero on their branches.\label{src-end-p032}\par

\begin{correction}\label{correction-3}
Read the preceding iteration as follows: iterate over $H$ for each coordinate
outside $H$, then charge $H$ at its original weights. Charging each later
coordinate at its weight after restriction gives an inequality in the wrong
direction. This is correction 1 in \path{CORRECTIONS.md}, reproduced below.
\end{correction}
% SOURCE P033, PROOF.md lines 139--144
\leavevmode\label{src-p033}Choose a count extremizer $f$ with $N=R_d$ and all local degrees $d$. Choose
one of its coordinates $i$ and let $H$ be the full relevant-coordinate set of
$A_i$; write $h=|H|$. Decompose $f=x_iA_i(x_H)+B(x_H,x_Y)$. On any control input
$\alpha$ with $A_i(\alpha)=\pm1$, both $x_i$ branches being Boolean forces $B=0$
for every $Y$. Thus $f_\alpha$ is a signed dictator and has $D=1/2$. The fraction
of such control assignments is exactly $p_i$.\label{src-end-p033}\par

% SOURCE P034, PROOF.md lines 146--152
\leavevmode\label{src-p034}Every degree-$d$ monomial meets $H$. For a monomial in $B$ disjoint from $H$, its
coefficient cannot change on restricting $H$, since every larger monomial
containing it would have degree above $d$. But that coefficient must vanish
on any active control assignment, since $B$ vanishes identically there. This
is a contradiction. A degree-$d$ monomial involving $i$ contains a nonempty
mask of $A_i$ (of size $d-1$), hence also meets $H$ for $d\ge2$. Consequently every
$f_\alpha$ has degree at most $d-1$.\label{src-end-p034}\par

% SOURCE P035, PROOF.md lines 154--155
\leavevmode\label{src-p035}Let $c\ge W_{d-1}$ with $c\ge1/2$. Formula (3), all original local degrees being $d$,
and the dictator branches imply\label{src-end-p035}\par

% SOURCE P036, PROOF.md lines 157--157
\begin{equation}
\frac{N}{2^d}\le\frac h{2^d}+\frac{p_i}{2}+(1-p_i)c.\tag*{}
\label{src-p036}
\end{equation}

% SOURCE P037, PROOF.md lines 159--159
\leavevmode\label{src-p037}Use (2) and $\lambda=Mp_i$ to obtain the essential inequality\label{src-end-p037}\par

% SOURCE P038, PROOF.md lines 161--162
\begin{equation}
\begin{aligned}N&\le F_d(c,\lambda)\\
&:=2^dc+\lfloor2(d-1)\lambda\rfloor+(1-2c)\lambda.\end{aligned}\tag{4}
\label{src-p038}
\end{equation}

% SOURCE P039, PROOF.md lines 164--164
\leavevmode\label{src-p039}Choose $i$ of minimum influence. Since $\sum_i p_i\le d$,\label{src-end-p039}\par

% SOURCE P040, PROOF.md lines 166--166
\begin{equation}
1\le\lambda\le\frac{dM}{N}.\tag{5}
\label{src-p040}
\end{equation}

% SOURCE P041, PROOF.md lines 168--169
\leavevmode\label{src-p041}Equations (4),(5) apply to the SAME chosen derivative. They do not assume
that every low-weight derivative has the same normal form.\label{src-end-p041}\par

% SOURCE P042, PROOF.md lines 171--171
\section{Exact energy grid and the higher low-weight spectrum}\label{src-p042}\label{src-end-p042}

% EXPANDED 6 November 2026: replaces the transcription of PROOF.md lines 173--197 (P043--P049)
\begin{editorial}
Revised 6 November 2026. In the first version, this section opened with a short paragraph
(frozen \path{PROOF.md}, lines 173--197). Steps 1--7 below write out every step of it.
The conclusions are unchanged: equation (6), and the fact that the
support of $A_i$ is a Reed--Muller codeword of normalized weight $\lambda$.
The paragraphs after Step 7 are reworded for clarity. They now treat the case $d=2$,
which the first version left implicit, and no mathematical claim is changed.
\end{editorial}
Section 3 bounds $N$ in terms of the real parameter $\lambda\ge1$ of one chosen
derivative. This section shows that $\lambda$ can take only special values.
Equation (6) puts $\lambda$ on the grid $(2/M)\mathbb Z$. The Reed--Muller weight
theorem then limits the values below $2$ to the finite list (8). Section 5
checks the values in (8) below $2$ and the necessary grid representatives (9) at
and above $2$.

\phantomsection\label{src-p043}\label{src-end-p043}\label{src-p044}\label{src-end-p044}\label{src-p045}\label{src-end-p045}%
\leavevmode Throughout this section, $f$ is a sign Boolean function of real degree at most $d$
with $d\ge2$, $i$ is a relevant coordinate, $k=d-1\ge1$, $M=2^k=2^{d-1}$ and
$\lambda=Mp_i$, as in Section 3. We use two coordinate systems on the cube: sign
coordinates $x_t\in\{+1,-1\}$ and bit coordinates $b_t\in\{0,1\}$, related by
\[
b_t=\frac{1-x_t}{2},\qquad x_t=1-2b_t ,
\]
so $x_t=+1$ corresponds to $b_t=0$. For a set $T$ of coordinates write
$x^T=\prod_{t\in T}x_t$ and $b^T=\prod_{t\in T}b_t$, with $x^\emptyset=b^\emptyset=1$.

\textbf{Step 1 (integer bit coefficients).} Every real function $u$ on $\{0,1\}^n$
has a unique multilinear expansion $u(b)=\sum_T c_T\,b^T$. For a set $U$, let
$1_U$ be the bit vector equal to 1 exactly on $U$. Then $b^T=1$ at $1_U$ if
$T\subseteq U$ and $b^T=0$ otherwise, so
\[
u(1_U)=\sum_{T\subseteq U}c_T\qquad\text{for every }U .
\]
M\"obius inversion on the subset lattice solves this triangular system:
\[
c_T=\sum_{U\subseteq T}(-1)^{|T|-|U|}\,u(1_U).
\]
(Substituting the first identity into the right side gives
$\sum_{V\subseteq T}c_V\sum_{V\subseteq U\subseteq T}(-1)^{|T|-|U|}$, and the inner sum
is $1$ if $V=T$ and $0$ otherwise.) Consequently, if $u$ takes only integer values,
every $c_T$ is an integer. This applies to the bit-output function
$g=(1-f)/2$, whose values are $0$ and $1$, and to $A_i$, whose values are $0,+1,-1$.

\textbf{Step 2 (the two degrees agree).} Substituting $b_t=(1-x_t)/2$ into $b^T$ gives
$\prod_{t\in T}(1-x_t)/2$, a multilinear polynomial in the $x_t$ of degree $|T|$.
So a multilinear bit polynomial of degree at most $e$ becomes a multilinear sign
polynomial of degree at most $e$. Substituting $x_t=1-2b_t$ shows the converse in the
same way. Multilinear expansions are unique in both coordinate systems, so the bit
degree and the real (Fourier) degree of any function on the cube are equal. In
particular $g=(1-f)/2$ has bit degree at most $d$, because its Fourier expansion is
$\tfrac12-\tfrac12 f$ and therefore has degree at most $d$.

\textbf{Step 3 (the Fourier grid of $f$).} Write $g=\sum_T c_T\,b^T$ with integer
$c_T$ (Step 1) and $c_T=0$ whenever $|T|>d$ (Step 2). For $S\subseteq T$, the coefficient
of $x^S$ in $b^T=2^{-|T|}\prod_{t\in T}(1-x_t)$ is $(-1)^{|S|}2^{-|T|}$. Since $f=1-2g$,
\[
\widehat f(S)=[S=\emptyset]-\sum_{T\supseteq S}(-1)^{|S|}\,c_T\,2^{1-|T|},
\]
where $[S=\emptyset]$ is $1$ for the empty set and $0$ otherwise. Each nonzero term has $|T|\le d$, so $2^{1-|T|}=2^{d-|T|}\cdot2^{1-d}$ is an integer
multiple of $2^{1-d}=1/M$. The constant $1=2^{d-1}\cdot 2^{1-d}$ is also a multiple of
$1/M$. Hence $M\widehat f(S)$ is an integer for every $S$.

\textbf{Step 4 (the grid of $A_i$).} Put the expansion $f=\sum_S\widehat f(S)x^S$ into
$A_i=\bigl(f(i{=}+1)-f(i{=}-1)\bigr)/2$. Terms with $i\notin S$ cancel. A term with
$i\in S$ becomes $\widehat f(S)\,x^{S\setminus\{i\}}$, since $(1-(-1))/2=1$. So
\[
\widehat{A_i}(S')=\widehat f\bigl(S'\cup\{i\}\bigr)\qquad(i\notin S'),
\]
and $\widehat{A_i}(S')=0$ when $i\in S'$. By Step 3, $z_{S'}=M\widehat{A_i}(S')$ is an
integer for every $S'$.

\textbf{Step 5 ($\lambda$ in terms of the $z_S$).} By Parseval,
\[
\lambda=M\,\mathbb E A_i^2=M\sum_S\widehat{A_i}(S)^2=\frac{\sum_S z_S^2}{M}.\tag*{}
\]
\phantomsection\label{src-p046}\label{src-end-p046}\label{src-p047}\label{src-end-p047}%
\textbf{Step 6 (parity).} At the all-$(+1)$ sign input every $x^S$ equals $1$, so
$A_i(+1,\ldots,+1)=\sum_S\widehat{A_i}(S)=\sum_S z_S/M$. The left side is $0$, $+1$ or
$-1$, so $\sum_S z_S=M\,A_i(+1,\ldots,+1)$ is an integer multiple of $M$. Because
$d\ge2$, $M=2^{d-1}$ is even, so $\sum_S z_S$ is even. For every integer $z$,
$z^2-z=z(z-1)$ is even, so $\sum_S z_S^2\equiv\sum_S z_S\equiv0\pmod 2$. Write
$\sum_S z_S^2=2m$ with $m$ an integer. Step 5 gives $\lambda=2m/M$, that is,
\begin{equation}
\lambda\in\frac2M\mathbb Z.\tag{6}
\label{src-p048}
\end{equation}
\phantomsection\label{src-end-p048}%

\phantomsection\label{src-p049}\label{src-end-p049}%
\textbf{Step 7 (the support of $A_i$ is a Reed--Muller codeword).} Regard $A_i$ as a
function on $n$ bit coordinates, adjoining unused coordinates if necessary so that
$n\ge k+2$. This leaves $p_i$, $\lambda$ and the degree bounds unchanged, because each
value is repeated equally often. Write $A_i$ in bit
coordinates, $A_i=\sum_T a_T\,b^T$. Its values are integers, so every $a_T$ is an integer
(Step 1). By Step 4, $\widehat{A_i}(S')$ is nonzero only when $\widehat f(S'\cup\{i\})$ is,
so the real degree of $A_i$ is $d_i-1\le k$. Hence $a_T=0$ for $|T|>k$ (Step 2).
Reduce each $a_T$ modulo two. The result $\bar A=\sum_T(a_T\bmod 2)\,b^T$ is a
polynomial over $\mathbb F_2$ of degree at most $k$. At any bit input, evaluating an
integer polynomial and then reducing modulo two gives the same result as reducing the
coefficients first and then evaluating. So $\bar A(b)\equiv A_i(b)\pmod 2$. Since
$0\equiv0$ and $\pm1\equiv1$, $\bar A(b)=1$ exactly when $A_i(b)\ne0$. Thus the
indicator of the support of $A_i$ is a word of the binary Reed--Muller code
$\operatorname{RM}(k,n)$. It is nonzero because $i$ is relevant. Its weight is
$w=2^n\Pr(A_i\ne0)=2^n\,\mathbb E A_i^2=2^np_i$, using $A_i^2=1$ on the support.
Relative to the minimum distance $2^{n-k}$ of $\operatorname{RM}(k,n)$, its
normalized weight is $w/2^{n-k}=2^kp_i=\lambda$. We reduce $A_i$ itself. Using only the real-degree bound $\deg(A_i^2)\le2k$ would give
membership in $\operatorname{RM}(\min\{2k,n\},n)$; reducing $A_i$ directly gives the
required order-$k$ bound.

% SOURCE P050, PROOF.md lines 199--201
\leavevmode\label{src-p050}We use the Kasami--Tokura theorem on Reed--Muller codewords of weight
less than twice the minimum distance (references below). For $k\ge2$, any nonzero
word of $\operatorname{RM}(k,n)$ with weight strictly below $2\cdot2^{n-k}$ has
normalized weight\label{src-end-p050}\par

% SOURCE P051, PROOF.md lines 203--203
\begin{equation}
\lambda=2-2^{-j},\quad\text{for an integer }j\ge0.\tag{7}
\label{src-p051}
\end{equation}

% SOURCE P052, PROOF.md lines 205--213
\leavevmode\label{src-p052}Since $w<2\cdot2^{n-k}$ is the same as $\lambda<2$, (7) applies exactly
when $\lambda<2$. The theorem also says which $j$ occur for given $n$ and $k$. We use
only the fact that $\lambda$ has the form (7), so we allow every $j\ge0$. As in Step 7,
we have $n\ge k+2$. The theorem applies to words of algebraic degree at most $k$. A
nonzero word of algebraic degree at most $k-1$ lies in $\operatorname{RM}(k-1,n)$, so
its weight is at least $2^{n-k+1}$, twice the minimum distance of
$\operatorname{RM}(k,n)$. The theorem is due to Kasami--Tokura (1970), with Berlekamp--Sloane, DOI
10.1109/TIT.1970.1054545. A published explicit statement and the two
normal forms appear in Carlet (2024), section 2, p.10613,
\url{https://www.aimspress.com/aimspress-data/math/2024/5/PDF/math-09-05-518.pdf} .\label{src-end-p052}\par

% SOURCE P053, PROOF.md lines 215--216
\leavevmode\label{src-p053}For $k\ge2$, combine (6) and (7). By (6), $M\lambda$ is an even integer, and by (7),
$M\lambda=M(2-2^{-j})=2^{k+1}-2^{k-j}$. For $j\le k-1$ this is even. For $j\ge k$ it is
either nonintegral ($j>k$) or odd ($j=k$). When $k=1$ ($d=2$), (6) says $\lambda$ is a
positive integer, so $\lambda<2$ forces $\lambda=1$, which is (8) below with $k=1$. (The
certificate in Section 5 starts at $d=4$, so this case is recorded only for completeness.)
Hence, for every $k\ge1$, every possible value of $\lambda$ below $2$ lies in the finite set\label{src-end-p053}\par

% SOURCE P054, PROOF.md lines 218--218
\begin{equation}
\{2-2^{-j}:j=0,\ldots,k-1\}.\tag{8}
\label{src-p054}
\end{equation}

% SOURCE P055, PROOF.md lines 220--226
\leavevmode\label{src-p055}The first terms are $1,3/2,7/4,15/8,31/16,\ldots$ . For $\lambda\ge2$ we use
only the grid (6). No further information about Reed--Muller weights in $[2,2.5)$
(Kasami--Tokura--Azumi) is used.\footnote{The attribution to third and higher below-two
weights is open. An AI referee's ablation finds that, given the grid (6), excluding
normalized weights in $(1,3/2)$ alone already gives the same rational, while the
ablation numbers in Section 5 disagree. The bound does not depend on that
attribution. See Correction 2, printed at the end of Section 5.}\label{src-end-p055}\par

% SOURCE P056, PROOF.md lines 228--228
\section{Finite scalar optimization, with the floor treated correctly}\label{src-p056}\label{src-end-p056}

% SOURCE P057, PROOF.md lines 230--233
\leavevmode\label{src-p057}Suppose $c$ bounds $W_{d-1}$. To exclude $N\ge n$, (5) puts $\lambda$ in $[1,t]$,
where $t=dM/n$. Since the set of possible $\lambda$ for $N\ge n$ is contained in
this interval, it suffices to show every allowed $\lambda$ in $[1,t]$ has
$F_d(c,\lambda)<n$. This excludes all $N\ge n$ at once, not merely $N=n$.\label{src-end-p057}\par

% SOURCE P058, PROOF.md lines 235--236
\leavevmode\label{src-p058}Use $c\in[1/2,(d-1)/2]$, checked inductively in the certificate. Within a
cell $h\le2k\lambda<h+1$, the value (4) is\label{src-end-p058}\par

% SOURCE P059, PROOF.md lines 238--238
\begin{equation}
2^dc+h+(1-2c)\lambda,\tag*{}
\label{src-p059}
\end{equation}

% SOURCE P060, PROOF.md lines 240--243
\leavevmode\label{src-p060}which is nonincreasing in $\lambda$. It is NOT safe to substitute $t$ directly
inside its floor. Below 2 we test every member of (8) not exceeding $t$.
For $\lambda\ge2$, the possible cell indices are $4k$ through $\lfloor2kt\rfloor$.
On the grid (6), the first value in cell $h$ is\label{src-end-p060}\par

% SOURCE P061, PROOF.md lines 245--245
\begin{equation}
\lambda_h=\frac2M\left\lceil\frac{hM}{4k}\right\rceil.\tag{9}
\label{src-p061}
\end{equation}

% SOURCE P062, PROOF.md lines 247--250
\leavevmode\label{src-p062}Discard it if $\lambda_h>t$ or $2k\lambda_h\ge h+1$. Every other grid value in
that cell has no larger budget. Thus (8),(9) form a complete finite cover
of all possible cases. Every such budget is strictly below $n$ in \path{caps.json}.
The resulting $c_d=(n-1)/2^d$ is valid by (4),(5) and the CHS identity.\label{src-end-p062}\par

% SOURCE P063, PROOF.md lines 252--254
\leavevmode\label{src-p063}Starting with $c_1=1/2,c_2=1,c_3=5/4$, proceed through $d=96$. The certificate
uses each already established $c_{d-1}$ before deriving $c_d$. There is no
circular use of the eventual constant. Some initial ceilings are:\label{src-end-p063}\par

% SOURCE P064, PROOF.md lines 256--266
\begin{center}
\begin{tabular}{rrr}
\toprule
$d$ & $R_d$ ceiling & $c_d$\\
\midrule
4 & 24 & $3/2$\\
5 & 54 & $27/16$\\
6 & 119 & $119/64$\\
7 & 254 & $127/64$\\
8 & 528 & $33/16$\\
9 & 1083 & $1083/512$\\
10 & 2199 & $2199/1024$\\
11 & 4439 & $4439/2048$\\
12 & 8928 & $279/128$\\
\bottomrule
\end{tabular}
\phantomsection\label{src-p064}\label{src-end-p064}
\end{center}

% SOURCE P065, PROOF.md lines 268--274
\leavevmode\label{src-p065}The frozen packet stores the most dangerous finite case at each first
excluded count, its complete case count, and a feasible scalar point at
the preceding count. Checker B independently enumerates all floor cells;
checker A uses an independently derived collapsed envelope. For the latter,
write $a=2c-1\in[0,k-1]$. Consecutive first grid values in occupied cells
$h<h'$ differ by at most $(h'-h)/(2k)+2/M$. Therefore their maximal budgets
increase by at least\label{src-end-p065}\par

% SOURCE P066, PROOF.md lines 276--277
\begin{equation}
\begin{aligned}(h'-h)\left(1-\frac a{2k}\right)-\frac{2a}{M}\\\ge1-\frac a{2k}-\frac{2a}{M}>0.\end{aligned}\tag*{}
\label{src-p066}
\end{equation}

% SOURCE P067, PROOF.md lines 279--284
\leavevmode\label{src-p067}The last strict inequality is checked exactly for every degree; universally
it follows for $k\ge3$ from $(k-1)/(2k)+2(k-1)/2^k<1$ (at $k=3$ this is $5/6$;
at $k\ge4$ the first term is below $1/2$ and the second is at most $3/8$).
Thus the rightmost occupied high cell suffices for checker A. It locates
that cell from the last grid point below $t$ and tests its first grid point.
Checker B uses none of this envelope shortcut.\label{src-end-p067}\par

% SOURCE P068, PROOF.md lines 286--292
\leavevmode\label{src-p068}Two additional certified variants quantify the inputs. Dropping the RM
gaps altogether and allowing a continuum $\lambda\ge1$ gives a bound below
2.2952887883. Retaining only $\lambda=1$ or $\lambda\ge3/2$ gives a bound below
2.2108998048. The third and higher below-two weights give the displayed
2.1974348818 bound. Each variant has its own inductive caps and both
checkers verify them. These are candidate, unrefereed consequences of the
same written lemmas, not claims about the exact Boolean maxima.\label{src-end-p068}\par

\begin{correction}\label{correction-5}
The attribution to third and higher below-two weights is open. The referee's
ablation finds that, given the energy grid, the single fact that no normalized
weight lies in $(1,3/2)$ already gives the same rational. Section 5's ablation
numbers disagree, and this has not been resolved; the bound does not depend
on that attribution. This is correction 2 in \path{CORRECTIONS.md}.
\end{correction}
% SOURCE P069, PROOF.md lines 294--294
\section{An elementary infinite tail}\label{src-p069}\label{src-end-p069}

% SOURCE P070, PROOF.md lines 296--297
\leavevmode\label{src-p070}We need a global bound, not just a horizon calculation. The CHS/Shinkar–Tal
lower family starts with $\Xi_1=x$ and, on disjoint variables, uses\label{src-end-p070}\par

% SOURCE P071, PROOF.md lines 299--299
\begin{equation}
\Xi_d=\frac{s+t}{2}\Xi_{d-1}(x)+\frac{s-t}{2}\Xi_{d-1}(y).\tag*{}
\label{src-p071}
\end{equation}

% SOURCE P072, PROOF.md lines 301--304
\leavevmode\label{src-p072}On a sign input exactly one summand is active, proving Booleanity by
induction. Its nonzero Fourier masks have size $d$. The two input copies
and the two fresh selectors remain relevant, so its count $L_d$ satisfies
$L_1=1$ and $L_d=2L_{d-1}+2$. Solving this recursion gives\label{src-end-p072}\par

% SOURCE P073, PROOF.md lines 306--306
\begin{equation}
L_d=3\cdot2^{d-1}-2=3M-2.\tag*{}
\label{src-p073}
\end{equation}

% SOURCE P074, PROOF.md lines 308--312
\leavevmode\label{src-p074}Thus $R_d\ge3M-2\ge2M$ for $d\ge2$. For the least-influence coordinate of its
count extremizer, (5) now yields $\lambda\le d/2$. Apply (4) with the ACTUAL
$c=W_{d-1}$, discard its floor downward, and divide by $2^d$. Since
$1/2\le W_{d-1}\le(d-1)/2$, the multiplier $2d-1-2W_{d-1}$ is positive and at
most $2d-2$. Hence\label{src-end-p074}\par

% SOURCE P075, PROOF.md lines 314--315
\begin{equation}
\begin{aligned}W_d&\le W_{d-1}+\frac{\lambda(2d-1-2W_{d-1})}{2^d}\\
&\le W_{d-1}+\frac{d(d-1)}{2^d}.\end{aligned}\tag{10}
\label{src-p075}
\end{equation}

% SOURCE P076, PROOF.md lines 317--319
\leavevmode\label{src-p076}No block-sensitivity theorem or inherited numerical block certificate is
required for (10). For $H\ge1$, write $d=H+j$ and sum $j\ge1$. The elementary
geometric-series identities are\label{src-end-p076}\par

% SOURCE P077, PROOF.md lines 321--321
\begin{equation}
\sum 2^{-j}=1,\qquad\sum j2^{-j}=2,\qquad\sum j^2 2^{-j}=6.\tag*{}
\label{src-p077}
\end{equation}

% SOURCE P078, PROOF.md lines 323--325
\leavevmode\label{src-p078}The first follows by summing the geometric series; applying $x\,d/dx$ to
$\sum_{j\ge1}x^j=x/(1-x)$ once and twice, then setting $x=1/2$, gives the
other two. Expansion of $(H+j)(H+j-1)$ gives\label{src-end-p078}\par

% SOURCE P079, PROOF.md lines 327--327
\begin{equation}
\sum_{d>H}\frac{d(d-1)}{2^d}=\frac{H^2+3H+4}{2^H}.\tag{11}
\label{src-p079}
\end{equation}

% SOURCE P080, PROOF.md lines 329--331
\leavevmode\label{src-p080}The computed caps are nondecreasing through $H$. Every smaller degree is
therefore bounded by $c_H$; for larger degrees telescope (10) and apply
(11). Consequently $C_{44}=\sup_d R_d/2^d$ satisfies\label{src-end-p080}\par

% SOURCE P081, PROOF.md lines 333--335
\begin{equation}
\begin{aligned}C_{44}&\le c_{96}+\frac{9508}{2^{96}}\\
&=\frac{174098727923678791596986402933}{79228162514264337593543950336}\\
&<2.1974348818.\end{aligned}\tag{12}
\label{src-p081}
\end{equation}

% SOURCE P082, PROOF.md lines 337--337
\leavevmode\label{src-p082}The exact improvement over round five's frozen rational is\label{src-end-p082}\par

% SOURCE P083, PROOF.md lines 339--340
\begin{equation}
\frac{76852852570035861934916069377}{79228162514264337593543950336}>\frac{1}{1000}.\tag*{}
\label{src-p083}
\end{equation}

% SOURCE P084, PROOF.md lines 342--343
\leavevmode\label{src-p084}This is an exact rational comparison. Rounding displayed decimals never
enters the proof.\label{src-end-p084}\par

% SOURCE P085, PROOF.md lines 345--345
\section{What the recurrence and the scalar limit do and do not establish}\label{src-p085}\label{src-end-p085}

% SOURCE P086, PROOF.md lines 347--347
\leavevmode\label{src-p086}The inherited valid block recurrence is\label{src-end-p086}\par

% SOURCE P087, PROOF.md lines 349--351
\begin{equation}
\begin{aligned}U(0,d)&=0,\\
U(b,d)&=\min\left(c_d,\frac d{2^d}+\max_{1\le j\le d}U(b-1,j)\right),\\
U(b,d)&=U(B_d,d)\quad\text{for }b>B_d.
\end{aligned}\tag*{}
\label{src-p087}
\end{equation}

% SOURCE P088, PROOF.md lines 353--360
\leavevmode\label{src-p088}Its proof and raw block duals are retained in \path{baseline/PROOF.md}, sections
5–6. In outline, restricting a maximum monomial of size $d$ costs $d/2^d$
in (3) and lowers block sensitivity by one: the nonzero full monomial
coefficient guarantees nonconstancy on its coordinates at every outside
input; this gives one sensitive block disjoint from all sensitive blocks
of the restricted function. Induct on $b$ and take the maximum over every
degree $j\le d$. Freezing a column above $B_d$ is valid because that class of
functions cannot grow further.\label{src-end-p088}\par

% SOURCE P089, PROOF.md lines 362--374
\leavevmode\label{src-p089}Checker B rechecks every retained Épi moment certificate and all 37
Fourier-lattice duals before re-running this recurrence with the inherited
$B_d$. Checker A separately computes integer recurrence rows using only
$B_d=d^2$, a consequence of the external Wellens theorem
$3\operatorname{bs}(\operatorname{bs}-1)\le2(d^4-d^2)$, Theorem 4.1, equation (4.3),
\url{https://arxiv.org/html/2005.00566} . Both reach exactly $c_d$ in EVERY final
column through degree 96. Thus the current block inputs are inactive at
the saturation rows of this run. The accepted global bound (12) has the
stronger elementary tail (11) and does not rely on those block premises.
The saturation thresholds and slack are in \path{analysis.json}. A smaller $B_d$
can help only if it crosses below its saturation threshold, with other
inputs fixed; a decrement that leaves it at or above that threshold has
zero effect in this recurrence.\label{src-end-p089}\par

% SOURCE P090, PROOF.md lines 376--381
\leavevmode\label{src-p090}There is also a precise limit for the scalar relaxation just specified.
Define its next count as the largest integer $N$ for which (4),(5),(6),(8)
are jointly feasible, allowing every high grid value $\ge2$. The packet's
feasible scalar points at its count ceilings and the strict exclusions at
the next integers prove that the computed counts are exactly optimal for
THIS relaxation. Those scalar points are not Boolean existence witnesses.\label{src-end-p090}\par

% SOURCE P091, PROOF.md lines 383--391
\leavevmode\label{src-p091}Let $c$ be the previous scalar cap. Its denominator divides $2^{d-1}$, so
$N_0=2^dc$ is an integer. The case $\lambda=1$ has budget
$N_0+2d-1-2c\ge N_0$ and satisfies $N_0\lambda\le dM$ whenever $c\le d/2$.
It belongs to every allowed low spectrum. Hence the next cap is at least $c$.
Every feasible scalar point at the largest count also has $N\ge3M-2$, so the
same estimates as in (10) bound the scalar increment by $d(d-1)/2^d$.
The maintained upper condition $c_d\le d/2$ continues after $d=96$ by induction,
since this increment is less than $1/2$ for $d\ge6$. Therefore the infinite
scalar sequence exists, is nondecreasing, and its limit lies in\label{src-end-p091}\par

% SOURCE P092, PROOF.md lines 393--394
\begin{equation}
\left[\frac{174098727923678791596986393425}{2^{96}},\ \frac{174098727923678791596986402933}{2^{96}}\right].\tag{13}
\label{src-p092}
\end{equation}

% SOURCE P093, PROOF.md lines 396--404
\leavevmode\label{src-p093}This pins its limit to $2.197434881723198\ldots$ with interval width
$9508/2^{96}<1.3\cdot10^{-25}$. It quantifies the end of this particular energy and
weight-spectrum relaxation. It is not an optimality theorem for all uses
of Reed–Muller structure. Full signed real-degree compatibility, sharper
control sets for actual higher-weight derivatives, correlations between
derivatives, or new block bounds crossing the saturation thresholds could
improve it. The known lower bound is $3/2$ from the family above, leaving a
candidate gap about 0.6974348817. The remaining gap cannot be removed by
merely iterating this already saturated scalar relaxation farther.\label{src-end-p093}\par

% SOURCE P094, PROOF.md lines 406--412
\leavevmode\label{src-p094}Related energy prior art: Potapov, \emph{On the number of relevant variables for
discrete functions}, Cryptography and Communications 17 (2025), 989–998,
\url{https://link.springer.com/article/10.1007/s12095-025-00797-4} , section 5,
also uses Fourier energy and changed edges to bound ternary relevance.
Its root-of-unity output encoding can double the real degree of a ternary
polynomial in our encoding. We credit the related method; no novelty is
claimed for the elementary real-energy argument in section 2.\label{src-end-p094}\par

\clearpage\section*{Corrections}\addcontentsline{toc}{section}{Corrections}\label{corrections}
\subsection*{Corrections to {\normalfont\texttt{PROOF.md}}}
\texttt{PROOF.md} is frozen so that its hash matches the receipts. A referee found two wording defects. Neither changes the mathematics.\par

1. Section 3, "Applying them successively to coordinates in H, or equivalently inducting on |H|". Read literally, this charges each later coordinate at its weight after restriction, and that inequality runs the wrong way. The correct reading: iterate over H for each coordinate outside H, then charge H at its original weights.\par

2. Sections 4 and 5 credit the third and higher below-two weights for the final bound. The referee's ablation finds that, given the energy grid, the single fact that no normalized weight lies in (1, 3/2) already gives the same rational. Section 5's ablation numbers disagree with that, and this hasn't been resolved. Treat the attribution as open. The bound doesn't depend on it.\par

External premises are R3 = 10 and the second Reed-Muller weight 3/2. The Kasami-Tokura-Azumi 1976 results aren't a premise. The claim in \texttt{KTA\_AUDIT.md} about a table omission is unverified.\par


\section*{External premises and cited results}
\addcontentsline{toc}{section}{External premises and cited results}
\label{premises}
The external premises identified in \path{CORRECTIONS.md} are $R_3=10$
and the second Reed--Muller weight $3/2$ (no normalized weight lies in
$(1,3/2)$). The degree-three result is cited through Tarannikov--Kirienko,
\url{https://eprint.iacr.org/2000/050}, and Krotov--Valyuzhenich,
\url{https://arxiv.org/html/2311.05566v2}, Section 3. These results are
cited, not reproved.

Section 4 cites the fuller below-two Kasami--Tokura spectrum, crediting
Berlekamp--Sloane for low binary weights: Kasami--Tokura (1970),
\url{https://doi.org/10.1109/TIT.1970.1054545}; see also Carlet (2024),
Section 2, p.~10613, linked there. Correction 2 leaves the need for the
higher below-two weights for this bound open.

Section 1 proves and cites the Chiarelli--Hatami--Saks identity. Section 7
cites Wellens's recurrence and block theorem, whose premises are not needed
for (12). Valyuzhenich's minimum-support classification is context, rather
than a premise of Section 2's energy inequality.

\textbf{KTA audit caveat.} The Kasami--Tokura--Azumi (1976) results are not
a premise. \path{KTA_AUDIT.md} reports that its original PDF endpoint and
the Elsevier API denied access, and that the full 1974 preprint was not
reconstructed. Its claim about an omission in the final table is
\emph{unverified}, as \path{CORRECTIONS.md} explicitly states. The proof
allows all energy-grid values at and above $2$ and uses no KTA exclusions
in the range $2$ to $2.5$. The audit's claimed higher-weight numerical gain
is also subject to correction 2.

\section*{Certificates and reproduction}
\addcontentsline{toc}{section}{Certificates and reproduction}
\label{certificates}
Pinned source folder:
\begin{center}\small
\url{https://github.com/rpinho/math-certificates/tree/09c84234137d3fadee1620dbc1c2ca58d2756ce2/c44/reed-muller-caps}
\end{center}
The frozen \path{PROOF.md} SHA-256, checked against \path{MANIFEST.json}, is
\begin{center}\small\path{af599ddff4732631a1c0565d5a5b3a90510d53955436b085b55883d85dcc4c47}\end{center}
The packet includes the following files and subfolders.
\begin{itemize}
\item Written sources: \path{PROOF.md}, \path{CORRECTIONS.md},
\path{README.md}, \path{PRIOR_ART.md}, \path{KTA_AUDIT.md}, and \path{review/}.
\item Exact inputs and generation: \path{caps.json}, \path{candidate.json},
\path{generate.py}, \path{analyze.py}, \path{analysis.json}, and \path{analysis.log}.
\item Checkers: \path{verify_caps_a.py} and \path{verify_caps_b.py}.
Checker B also reads \path{baseline/epi_certificate.json} and
\path{baseline/spectral_certificate_round3.json}; \path{baseline/} retains
the inherited proofs and certificates.
\item Receipts and hashes: \path{check_a.log}, \path{check_b.log},
\path{verification_a.json}, \path{verification_b.json},
\path{FINAL_RECEIPT.json}, \path{MANIFEST.json}, \path{SHA256SUMS},
and \path{PROCESS_LEDGER.json}.
\item Finite controls: \path{control_boolean.cpp}, \path{control_ternary.cpp},
\path{control_corruption.py}, \path{control_corruption.json},
\path{boolean.log}, \path{ternary.log}, and \path{corruption.log}.
\path{jobs/} retains the original reproduction wrappers.
\end{itemize}
From a writable copy of the pinned folder, using Python 3.10 or later
(standard library), the two checker commands are
\begin{verbatim}
nice -n 10 python3 -B verify_caps_a.py candidate.json
nice -n 10 python3 -B verify_caps_b.py candidate.json
\end{verbatim}
The scripts write \path{verification_a.json} and \path{verification_b.json}
beside themselves. Both were replayed for this PDF and returned PASS with
the rational in (12). They check arithmetic conditional on the written
lemmas and external premises, not those universal statements themselves.

\section*{Credit and later work}
\addcontentsline{toc}{section}{Credit and later work}
\label{credit}
Following the source README and prior-art note, credit goes to Chiarelli,
Hatami and Saks for the weighted framework (including the corrected
duplication in Chiarelli's thesis); Wellens for the recurrence; Berlekamp
and Sloane, Kasami, Tokura and Azumi for the weight spectra; Valyuzhenich
for minimum-influence structure; BakeLab/Épi for inherited exact
polynomial certificates; and Shinkar and Tal for the lower family also
credited in Section 6. The degree-three citations are due to Tarannikov
and Kirienko, and Krotov and Valyuzhenich. Potapov's related ternary
spectral-energy method is credited in Section 7. The source prior-art
note and KTA audit also discuss Carlet, Rolland, Klein and Storme, Krotov's
public discussion, Lou and Wang, and Wei et al.; these contextual
references do not supply additional premises for (12). No novelty,
priority, or optimality is claimed.

Later work in the same repository (\path{small-degree-classification})
gives $C_{44}<1.9916994577$, also unrefereed; this PDF concerns the requested
$2.1974348818$ proof.
\end{document}
